\documentclass[10pt,a4paper]{article}
\usepackage{fontspec,amsmath,amssymb,graphicx}
\usepackage[margin=18mm]{geometry}
\setmainfont{Times New Roman}
\setlength{\parindent}{0pt}
\setlength{\parskip}{5pt}
\allowdisplaybreaks
\sloppy
\emergencystretch=2em
\begin{document}
\section*{BT\_5\_4\_3\_14 — Lời giải bổ sung}
Đặt $v(x,y)=u(x,y)+u(x,-y)$. Đây là hàm điều hòa chẵn theo $y$, và $v(x,0)=0$, $v_y(x,0)=0$. Từ $v_{yy}=-\Delta_xv$, tất cả đạo hàm theo $y$ trên siêu phẳng đều 0 theo quy nạp. Hàm điều hòa giải tích, nên $v=0$ gần siêu phẳng, rồi trên toàn miền nhờ tính duy nhất giải tích. Suy $u(x,-y)=-u(x,y)$.\end{document}